% Generated from completed comparison evidence.
\section{Conclusiones de comparación y fidelidad}
El banco de 1050 ajustes conserva diferencias entre estimadores antes y después de añadir restricciones. En 102 de 105 observaciones no existe solución admisible Raw; con H+B quedan 0 observaciones sin solución. El resultado se interpreta como cobertura de configuraciones fijas bajo los umbrales, no como selección automática de un modelo en una medición cuyo limpio no está disponible.

El contraste Raw fijo PCP-TLS produce una mediana de diferencia de $e_\omega$ de -203.2 pp bajo contaminación impulsiva y -83.14 pp bajo gaussiano IID. El signo negativo favorece PCP; los diagnósticos finitos incluyen los estados de convergencia documentados y la aceptación se comprueba aparte. Esta comparación conserva la misma observación por par y no transforma una ventaja condicionada en superioridad universal.

El criterio modifica la preferencia: con gaussiano IID, la diferencia de $e_u$ PCP-TLS es 2.981 pp, aun cuando PCP tiene menor $e_\omega$. Con ruido espacial correlacionado la diferencia de $e_\omega$ cambia a 16.91 pp. Un error menor que el de otro estimador puede seguir fuera de tolerancia; los conteos de cobertura mantienen esa distinción.

Los contrastes pareados y el control identity con el mismo prior impiden atribuir toda limpieza espacial al modelo dinámico. Un error bajo en velocidad puede coexistir con vorticidad o cierres constitutivos degradados. En fluidos de memoria, el esfuerzo no se identifica mediante reconstrucción de velocidad sola; en viscoplasticidad se evalúa bajo la hipótesis explícita de ley material conocida.

El filtro k\_c=10 impone un piso de 10.84~\% de error relativo de vorticidad en Multiescala NS contra su referencia completa. Por ortogonalidad espectral, una salida limitada a esa banda no puede alcanzar la tolerancia del 10~\%, incluso con observación perfecta. Aumentar aceptación con una puerta del 25~\% requiere reconocer esa pérdida de señal, además del ruido que se elimina.

La suficiencia de simulación depende del observable, la malla y el intervalo. Las referencias finas del mismo solver aportan convergencia numérica, mientras las soluciones analíticas y controles temporales examinan causas diferentes. Los errores entre mallas no acreditan validación experimental, y el control viscoplástico corto conserva su horizonte.

La sensibilidad a rango, intensidad y tolerancia explica límites de las configuraciones; no utiliza el limpio para escoger el mejor hiperparámetro desplegable. La contribución comprobable es una comparación trazable de observación, estimador, prior físico, referencia y recursos, con resultados negativos preservados.
